\section{Related Works}
\paragraph{Feed Forward Reconstruction}
% Quickly Introduce VGGT and other 3D reconstruction methods
Traditional multi-view reconstruction pipelines \cite{Lowe:2004:DIF:993451.996342, Hartley2003, schoenberger2016sfm} rely on sequential stages of feature extraction, matching, and geometric optimization, making them computationally expensive and sensitive to initialization. Feed-forward reconstruction models address this by directly regressing 3D structure and camera parameters from images in a single forward pass.
Feed-forward reconstruction models address this by directly regressing 3D structure and camera parameters from images in a single forward pass. DUSt3R \cite{wang2024dust3rgeometric3dvision}and MASt3R \cite{leroy2024groundingimagematching3d} established this paradigm by predicting dense point maps and relative camera poses from image pairs without requiring known intrinsics, relying on post-hoc global alignment to extend beyond pairs. Methods like VGGT \cite{wang2025vggtvisualgeometrygrounded} or MapAnything \cite{keetha2026mapanythinguniversalfeedforwardmetric} advanced this further by processing up to hundreds of images simultaneously in a single forward pass, jointly predicting camera parameters, depth maps, point maps, and 3D point tracks.
\paragraph{Bundle Adjustment}
% Explain quick idea behind it
Bundle adjustment (BA) is the gold standard for refining camera poses and 3D structure in multi-view reconstruction. Formally introduced by Triggs et al. \cite{10.1007/3-540-44480-7_21}, it minimizes the total reprojection error over all cameras and scene points simultaneously via the Levenberg-Marquardt algorithm, exploiting the sparse block structure of the Jacobian to keep the problem tractable. In classical SfM pipelines such as COLMAP \cite{schoenberger2016sfm}, BA is applied both locally after each image registration and globally as a final refinement, preventing drift and enforcing geometric consistency across the full reconstruction. While highly effective given sufficient image overlap and reliable correspondences, classical BA is sensitive to initialization and provides no mechanism to recover from fundamentally incorrect pose estimates. This limitation becomes particular acute in the sparse-view setting.
\paragraph{Superquadrics}
% What are superquadrics? How can they be useful for 3D scene representation?
Superquadrics are a family of parametric geometric primitives, first introduced by Barr \cite{1673799} and later popularized for 3D shape recovery by Solina and Bajcsy \cite{article}, capable of representing a wide range of shapes — from boxes to cylinders to ellipsoids — with a small set of parameters encoding position, orientation, scale, and two shape exponents. Their compactness and analytical differentiability make them attractive as structural scene representations. SuperDec \cite{fedele2026superdec3dscenedecomposition} builds on this foundation by presenting an approach for creating compact 3D scene representations through decomposition into superquadric primitives, designing a novel architecture that efficiently decomposes point clouds of arbitrary objects into a compact set of superquadrics. Rather than targeting photorealistic reconstruction, SuperDec solves the decomposition locally on individual objects and leverages instance segmentation to scale to full 3D scenes. In our work, we use SuperDec to obtain a compact geometric abstraction of the scene from the ground truth point cloud, which we then use as a structural prior to constrain and guide bundle adjustment optimization.